% =====================================================================
% capa.tex -- pagina de rosto / capa, inspirada na diagramacao de
% documentos do NRC (NUREG-1537) e do AP1000 DCD, sem reproduzir
% brasoes, selos ou logotipos oficiais.
% =====================================================================
\clearpage
\newgeometry{top=1.6cm,bottom=2.5cm,left=3.0cm,right=2.3cm}
\thispagestyle{empty}
\singlespacing
\begin{center}

\rule{\textwidth}{1.6pt}\\[2pt]
\rule{\textwidth}{0.6pt}\\[14pt]

{\sffamily\large COMISSÃO NACIONAL DE ENERGIA NUCLEAR}\\[2pt]
{\sffamily\small CNEN/RMB}

\vspace{0.8cm}

{\sffamily\bfseries\LARGE RELAT\'ORIO PRELIMINAR\\[3pt] DE AN\'ALISE DE SEGURAN\c{C}A}\\[8pt]
{\sffamily\normalsize (RPAS)}\\[3pt]
{\sffamily\normalsize Requerimento de Permiss\~ao de Constru\c{c}\~ao}

\vspace{0.6cm}

{\sffamily\bfseries\Large CAP\'ITULO 1}\\[4pt]
{\sffamily\large A INSTALA\c{C}\~AO}

\vspace{0.35cm}
{\sffamily\small PROPOSTA -- MINUTA DE TRABALHO PARA DISCUSS\~AO INTERNA}

\vspace{0.5cm}

\begin{tabular}{@{}>{\raggedright}p{0.32\textwidth}p{0.56\textwidth}@{}}
\textbf{Identifica\c{c}\~ao do documento:} & RMB-N00-00-LN-G-7240-RA-001 \\
\textbf{Revis\~ao:} & 0 (minuta) \\
\textbf{Data:} & Setembro de 2026 \\
\textbf{Idioma:} & Portugu\^es \\
\textbf{Refer\^encia de formato:} & USNRC NUREG-1537, Partes 1 e 2 \citep{nureg1537p1,nureg1537p2}\\
\end{tabular}

\vfill

\begin{minipage}{0.9\textwidth}
\centering\small
\textit{\textbf{AVISO: }Este documento \'e uma proposta preliminar de estrutura e conte\'udo para o
Cap\'itulo 1 do RPAS, elaborada pela FERMIUM TECNOLOGIA NUCLEAR com base no formato sugerido
pela NUREG-1537 (aplic\'avel a reatores n\~ao-de-pot\^encia / multi\-pro\-p\'osito) e na
apresenta\c{c}\~ao gr\'afica de Relat\'orios Finais de An\'alise de Seguran\c{c}a (FSAR) de
projetos certificados.}
\end{minipage}

\vspace{0.5cm}
\rule{\textwidth}{0.6pt}\\[2pt]
\rule{\textwidth}{1.6pt}

\end{center}
\restoregeometry
\clearpage
